\documentclass[11pt]{article}
\usepackage[margin=1in]{geometry}
\usepackage{amsmath,amssymb,booktabs}
\usepackage[hidelinks]{hyperref}
\newcommand{\fl}{\operatorname{fl}}
\title{A dyadic choice of constants for the seed cover of\\ \emph{The maximum number of
mutually unbiased bases in dimension six}}
\author{}
\date{}
\begin{document}
\maketitle

\noindent
The preprint fixes $N=112$, $w=1/112$, $R=.032$, $r_s=.0320001$ and
$\mathrm{RN2}=(\fl(.032\cdot112))^2-10^{-4}$ for its seed cover, and then spends most of
Appendix~A.3 justifying that these numbers survive binary64. This note replaces them by
\[
  N=128=2^7,\qquad w=2^{-7},\qquad R=2^{-5}=\tfrac1{32},\qquad r_s=R,\qquad
  \mathrm{RN2}=(NR)^2=16,
\]
for which every quantity the seed stage compares is an exact binary64 number. The symbolic
lemmas of \S2 are unchanged; what disappears is the floating-point analysis, the $+10^{-7}$
enlargement of the seed ball, and the code-generation premise about \texttt{pow}.

\section{The choice is forced}

Three of the paper's own premises bound the pair $(N,R)$:
\begin{enumerate}
\item[(P1)] bin differences are stored in \texttt{int8\_t} and \S7.3 asserts they lie in
  $[-111,111]$; in general they lie in $[-(N-1),N-1]$, so $N-1\le 127$, i.e. $N\le 128$;
\item[(P2)] the field bound $\mathrm{Rs}=2\pi R<.203$ of \S A.13 gives $R<.032309$;
\item[(P3)] the \emph{cheap} reuse test (2.18) can fire only if the ball of radius $R$ about one
  midpoint contains the whole box of another, i.e.\ only if
  $R\ge\max_{|h|\le\frac12}\|k+h\|/N$ for some $k\in\mathbb Z^9\setminus\{0\}$. The minimum is at
  $k=e_1$ and equals $\sqrt{(3/2)^2+8\cdot(1/2)^2}/N=\sqrt{17}/(2N)=2.0616/N$; the
  triangle-inequality form $1/N+3/(2N)=2.5/N$ is a strict overestimate. The \emph{dual} test (2.19)
  needs to cover only the feasible portion of a bin, a 5-dimensional slice of the box, so it can
  merge below this floor: (P3) bounds the cheap path and does not prove that reuse dies.
\end{enumerate}
Now impose exactness. A bin midpoint is $(b+\tfrac12)/N=(2b+1)/(2N)$, exactly representable in
binary64 if and only if $2N$ is a power of two, i.e.\ if and only if $N$ is. By (P1) the only powers
of two at most $128$ are $64$ and $128$, and $N=64$ is excluded by cost rather than by an
inequality: with $R\le.03231$ its effective separation $R-3/(2N)\le.0087$ is half the bin spacing
$1/64$, so no bin's own midpoint lies within the cheap test's reach and the entire merging load
falls on the dual test. $N=128$ works: $NR=4$, hence $(\mathrm{RN})^2=16$ exactly, and
$\delta=5/256$, i.e.\ $2.5$ bin spacings. For $R$ it suffices that $R$ be dyadic; in the window
$[2.0616/128,\,.032309]$ the dyadic values are $m/256$, $m=5,6,7,8$. The largest, $R=1/32=8/256$,
retains fewest centres and is the only one giving $NR=4$, a power of two. $112$ is not selectable at
all (midpoint denominator $224$), $96$ neither ($192$).

One caveat: $R=1/32$ is not the largest \emph{dyadic} admissible. $R=33/1024=.032227$ is dyadic,
keeps $NR=4.125$ and $(NR)^2=17.015625$ exact, and would (on a packing model) retain some $16\%$
fewer centres --- at the price of leaving $\mathrm{Rs}=2\pi R=.202485$ only $0.25\%$ below $.203$
instead of $3.3\%$, which is exactly the near-tight constant the exercise is removing.

\section{The constants and the inequalities they must satisfy}

\begin{center}
\begin{tabular}{llll}
\toprule
quantity & paper & dyadic & note\\
\midrule
$N$ & $112$ & $128$ & largest allowed by (P1)\\
$w$ & $1/112$ (stored $\approx$) & $2^{-7}$ & exact\\
$R$ & $.032$ & $2^{-5}$ & $NR=4$ exact\\
$r_s$ & $.0320001$ & $1/32$ & no enlargement\\
$\mathrm{RN2}$ & $(\fl(.032\cdot112))^2-10^{-4}$ & $16$ & exact, no slack\\
$3/(2N)$ & $3/224=.013393$ & $3/256=.011719$ & bin circumradius\\
$\delta=R-3/(2N)$ & $.018607$ & $5/256=.019531$ & effective separation\\
$\mathrm{Rs}=2\pi R$ & $.201062$ & $\pi/16=.196350$ & field bound\\
\bottomrule
\end{tabular}
\end{center}

\begin{center}
\begin{tabular}{lcccc}
\toprule
inequality & paper & slack & dyadic & slack\\
\midrule
$3/(2N)\le R$ & $.013393\le.032$ & $2.39\times$ & $.011719\le.03125$ & $2.67\times$\\
$R<.203/(2\pi)$ & $.032<.032309$ & $0.96\%$ & $.03125<.032309$ & $3.3\%$\\
$R\ge\sqrt{17}/(2N)$ & $.032\ge.018407$ & $1.74\times$ & $.03125\ge.016106$ & $1.94\times$\\
$N-1\le127$ & $111$ & $16$ & $127$ & $0$\\
$\binom{N+4}{5}<2^{31}$ & $1.604\cdot10^8$ & $13.4\times$ & $3.093\cdot10^8$ & $6.9\times$\\
$NR$ exact & $448/125$, no & --- & $4$, yes & ---\\
$(NR)^2$ exact & $12.844956$, no & --- & $16$, yes & ---\\
slack below $(NR)^2$ & $10^{-4}$ & --- & $0$ & ---\\
\bottomrule
\end{tabular}
\end{center}

\noindent
Two slacks get worse, and both are integer type bounds rather than numerical ones: the
\texttt{int8\_t} range is now exactly full, and the 32-bit tuple count halves its margin. The first
is a \emph{soundness} margin, not a cost margin: \texttt{covout} stores $A.b_j-p.b_j$ in an
\texttt{int8\_t}, and a wrapped difference would understate the squared distance and could certify a
coverage that does not hold. So $N\le128$ must be stated as an invariant with that reason, and
$N=128$ is the last admissible value; any finer cover has to widen the type first.

\section{What leaves the proof}

\begin{enumerate}
\item The enlargement $r_s=.0320001$ (five occurrences, and \texttt{atlas.cpp}) becomes $1/32$. The
  closing paragraph of A.3 --- the nine division-rounded seed phases differ from their rational
  midpoints by less than $10^{-14}$, absorbed by the enlarged radius --- is deleted: the division
  is exact.
\item A.3's comparison radius $\mathrm{RN2}=\fl((\fl(.032\cdot112))^2)-.0001$ becomes
  $\mathrm{RN2}=(NR)^2=16$: with $R\cdot N=4$ exactly, $\texttt{pow}(4,2)=16$ both under a
  correctly rounded \texttt{pow} and under the multiply lowering. The paragraph of \S7.1 that
  inspects the saved disassembly and has to be repeated for every new build is replaced by one
  sentence. This is not cosmetic: a rebuild of \texttt{covout.cpp} with this machine's g++~16.2 and
  the paper's own flags emits two \texttt{call pow@plt}, so the recorded gcc~13.3 binary satisfies
  \S7.1 and a rebuild here does not; at $(128,1/32)$ the condition is vacuous. (The second
  \texttt{pow}, inside $T=(\sqrt{\mathrm{RN2}}+3/2)^2$, has a non-exact argument, but $T$ only
  selects which centres are \emph{tried}.)

  The $-10^{-4}$ is deleted rather than replaced, and this is the one place where dropping a margin
  is safe. It covered the \emph{threshold's} own computation error ($<10^{-12}$ at $N=112$); at
  $N=128$ that error is zero. Both tests stay conservative against the exact threshold: the cheap
  test evaluates $s+a+9/4$, exactly the supremum of $\|k+h\|^2$ over the box, and the dual test
  evaluates $s+9/4+\texttt{supdual}$, where the returned value already carries $+.01$ over a bound
  accurate to $10^{-4}$, so it exceeds $2k\cdot h$ by about $0.009$ --- three orders above the single
  binary64 rounding of the sum ($\le1.8\cdot10^{-15}$). Hence ``computed $<16$'' implies
  $\|k+h\|^2<16$ on both branches.

  With the threshold exact, the seed arithmetic collapses to one integer comparison. The cheap test
  reads $s+a+9/4<16$ with $s=\|k_X\|^2$ and $a=\|k_Y\|_1$ non-negative integers, i.e.\
  \[
    s+a\le 13,
  \]
  with no floating point in the decision at all. That single line is where the whole of A.3's seed
  arithmetic now lives, and it is why the choice among $16$, $16-10^{-4}$ and $16-2^{-14}$ is
  cosmetic: the left side is quantised to integers plus $9/4$, so all three give the same cut.
\item A.3's opening ``$n=112$ and the stored $w$ approximates $1/n$'' becomes ``$n=128$ and
  $w=1/128$ exactly''. The gap comparisons are then exact except for the one-sided $10^{-8}$
  inflation that keeps extra tuples; the $10^{-14}$ error bound becomes zero.
\end{enumerate}

\noindent
What does not leave: the $2\cdot10^{-14}$ evaluation error of \texttt{expair} (\S A.2) and the
\texttt{sumok}/\texttt{extend}/\texttt{supdual} margins. Those concern the polynomial evaluation
of $e(x)$ and are independent of $N$; they keep their form with $N=128$ substituted.

\section{The re-derivation, in the paper's notation}

The statements of Lemmas 2.4--2.9 are uniform in $N$ with $w=1/N$ and are used verbatim. The
three numeric steps become:
\begin{itemize}
\item \S2.2: $N=128$, $w=1/N=2^{-7}$; the gap thresholds $1-(l-1)/N$ and $1-(l+2)/N$ are exact
  multiples of $2^{-7}$.
\item \S2.4: a retained midpoint covers its bin with radius
  $\sqrt{9}/(2N)=3/256<1/32$; a newly inserted midpoint needs radius
  $1.5/128=.0117<.03125$; the reuse test succeeds only when the bound of Lemma 2.9 is below
  $\mathrm{RN2}=(NR)^2=16$, with no slack needed (see \S3). The integer and phase parts of that
  bound are exact; the dual term keeps the outward margin it already has in A.3.
\item Prop.\ 2.10: the written centres are exact binary64 numbers, so the radius is $1/32$ with no
  enlargement.
\end{itemize}

\noindent
\emph{The seed that must never be rejected.} The exact Fourier matrix $e(ij/6)$ has $X$ bins
$(21,42,64,85,106)$ with $s=2$ and $Y$ bins $(64,21,42,85,106)$ at $N=128$. It passes the
orientation filter ($21\le106-85+1=22$), both gap filters ($g=.171875$ against $.179688$ and
$.156250$), the shared-interval test, and \texttt{sumok} in all six directions, with thinnest
margin $3.10\cdot10^{-5}$ (against $4.19\cdot10^{-5}$ at $N=112$: the test tightens as $w$ shrinks,
and still sits six orders above the $10^{-8}$ additions). Its shared displacement is exactly
$h_s=-\tfrac12$ at every $N$, so the critical seed lies on the closed boundary of the box
$[-\tfrac12,\tfrac12]^9$: every inequality that uses that box --- (2.12), (2.17), (2.18) and the
initial interval of \texttt{extend} --- must stay inclusive.

\section{The measured cost}

The authors' \texttt{covout.cpp} and \texttt{seed.cpp} (Apache-2.0, the versions used for the
recorded run) were rebuilt with the three constants changed and run to completion on the same
machine.

\begin{center}
\begin{tabular}{lrrr}
\toprule
 & $(112,.032)$ & $(128,1/32)$ & ratio\\
\midrule
X tuples surviving the filters & $22{,}829$ & $32{,}018$ & $1.40$\\
pairs tested & $151{,}744{,}086$ & $267{,}628{,}436$ & $1.76$\\
retained centres & $249{,}878$ & $229{,}979$ & $0.92$\\
cover CPU time & $402$\,s & $517$\,s & $1.29$\\
$\binom{N+4}{5}$ & $160{,}389{,}488$ & $309{,}319{,}296$ & $1.93$\\
$N\binom{N+3}{4}$ & $774{,}294{,}080$ & $1{,}499{,}729{,}920$ & $1.94$\\
\bottomrule
\end{tabular}
\end{center}

\noindent
The pair count cannot be extrapolated linearly from the progress lines. Pairs added per $X$ tuple,
from the checkpoints: at $N=128$, $3{,}845$ over the first $10{,}000$, $7{,}903$ over the second,
$11{,}975$ over the third and $15{,}062$ over the last $2{,}018$; at $N=112$, $3{,}749$ / $8{,}238$ /
$11{,}266$. The density grows roughly in proportion to the $X$ index, because the walk is
lexicographic and a larger largest bin admits more $Y$ tails, so the total scales like $T^2/2$ in the
index and a two-checkpoint extrapolation undershoots by about $1.3\times$; a linear-density fit on
the first two blocks reproduces both totals to better than $0.6\%$. The pair growth is accounted for
rather than fitted: an exact count of every filter but the shared-interval test and the $Y$-side
\texttt{sumok} gives $2.3706\times$, the residual acceptance rate falls to $0.744$, and
$2.3706\times0.744=1.7637$.

The change makes the computation exact but not cheaper. The cover costs $29\%$ more CPU; the centre
count falls $8\%$; at Stage A's recorded $8.08$\,ms per centre that saves about $160$\,s against the
cover's extra $115$\,s, i.e.\ a wash on the $4{,}470$\,s run. A packing model calibrated at 112
predicted $0.65$--$0.78\times$ the centres and $0.92\times$ was measured: the seed set is a lattice,
not a continuum, so do not price the change on the model. The output file is well formed
(\texttt{128 229979}, every bin in $[0,127]$, i.e.\ the \texttt{int8} range used to its last unit).

\section{Three more things worth fixing in the same pass}

\paragraph{Stage B, $k\le10$.} The base grid is $n=2L+1$ with $L=\lceil E/w+.500001\rceil$ and
$E\approx R_{\text{parent}}$, so the box cap $n^k\le5\cdot10^7$ gives $k\le5$ in the first two
passes ($n=23$ at $(.07,.013)$ and $n=31$ at $(.016,.005)$) and $k\le7$ in the third ($n=11$ at
$(.011,.0038)$). The paper's error budgets assume $k\le10$, which is the safe direction, but the
stated limit should read $k\le5,5,7$.

\paragraph{Stage C, the cluster cutoff.} \texttt{cluster()} runs once, after the four bin layers, so
the $M$ in $(.042-1.11/M)^2$ is always the deepest grid $96$; the paper never says which $M$.
$.042$ is the clustering's \emph{output} ball radius, and the threshold should be $(R-X_M)^2$ with
$X_M=\sqrt{29/6}/(2M)=.01145$ (the code's own \texttt{xnorm}), not $1.11/M$ --- $1.11$ exceeds
$\sqrt{29/6}/2=1.09924$ by $1\%$. The defect is only the squaring, which loses the sign: the base is
negative for $M\le26$, the merge radius is non-monotone, and on a one-layer invocation the implied
radius bound is $.1421>.1$, at which every record fails. The fix is one line, $\tau=\max(0,R-X_M)$,
and it yields the clean invariant $s\le R+1.001\cdot10^{-9}$ instead of $s\le R-1.12\cdot10^{-4}$.
Keep $.042$: measured on the pipeline's own records, $.03$ multiplies the ball count by $5.7\times$
($306\to1755$) and the stage's CPU by $4\times$ and still leaves $12/12$ hardest records over the
graph cap, while $.06/.075/.09$ shrink the balls but cost the $U$-test its discriminating power.

\paragraph{Stage C, the graph cap.} The real defect in \S6 is the $5\cdot10^6$ step cap: all $163$
records that exhausted it terminate \emph{negative} at $5\cdot10^7$ (median $9.9\cdot10^6$ steps,
maximum $39{,}038{,}705$, zero candidates). Raising the cap to $10^8$ costs about $1.9\cdot10^9$
steps, some $48$\,s of CPU, and deletes the fallback round --- five rows of Table~1, the bad-file
list, and 5 of the 100 recorded invocations.

\end{document}
